\subsection{ArUco: словарь и детекция}
\label{sec:python-aruco}

Для роботов существуют специальные визуальные метки (маркеры), которые легко распознать на изображении. Самые популярные — \textbf{ArUco} и \textbf{AprilTag}.

\subsubsection*{Что такое ArUco-маркер?}
Это черно-белый квадрат с уникальным узором внутри. В узор закодирован \textbf{ID} --- номер рисунка. ID только называет метку: ни расстояния, ни координаты из числа ID не получить. Что означает метка с данным номером («цель», «северо-запад», «посадочная площадка»), задаёт договорённость --- таблица в программе.
\begin{itemize}
    \item \textbf{Словарь (Dictionary):} набор допустимых рисунков, \verb|DICT_<сетка>_<сколько номеров>|. С Пионером используют маркеры $6\times6$: учебные метки симулятора и карты --- \texttt{DICT\_6X6\_250} (250 маркеров), часть старых примеров этой главы --- \texttt{DICT\_6X6\_50}. Детектор находит только метки своего словаря: рисунок другого словаря виден, но не распознаётся.
    \item \textbf{AprilTag:} Другой популярный формат, поддерживаемый той же библиотекой (например, \texttt{DICT\_APRILTAG\_36h11}).
\end{itemize}

\begin{warnbox}[Словарь в симуляторе --- точно DICT\_6X6\_250]
В настоящем OpenCV первые 50 маркеров \texttt{DICT\_6X6\_250} совпадают с \texttt{DICT\_6X6\_50}, поэтому на дроне метки 0--49 находят оба словаря. Симулятор сравнивает словари строго: детектор с \texttt{DICT\_6X6\_50} или \texttt{DICT\_4X4\_50} не найдёт на сцене ничего. Пишите \texttt{DICT\_6X6\_250}.
\end{warnbox}

\subsubsection*{Детекция шаг за шагом}
Для работы с камерой и маркерами мы используем библиотеку \textbf{OpenCV} (\texttt{cv2}). Детектор создаётся один раз, а кадров обрабатывает сколько угодно:

\begin{codebox}[Детектор]
\begin{verbatim}
dictionary = cv2.aruco.getPredefinedDictionary(cv2.aruco.DICT_6X6_250)
detector = cv2.aruco.ArucoDetector(dictionary)
corners, ids, rejected = detector.detectMarkers(frame)
\end{verbatim}
\end{codebox}

\verb|detectMarkers| ищет на кадре квадраты, сверяет их узор со словарём и возвращает три значения:
\begin{itemize}
    \item \verb|ids| --- номера найденных маркеров: массив строк, в каждой одно число, поэтому номер берут как \verb|int(ids[i][0])|. \textbf{Если меток нет, \texttt{ids} равен \texttt{None}}, а не пустому списку. Это штатный исход: кадр есть, меток в нём нет. Его нужно отличать от отсутствия кадра (\verb|frame is None|) из прошлого урока.
    \item \verb|corners| --- углы маркеров: \verb|corners[i][0]| --- четыре точки $(x, y)$ маркера с номером \verb|ids[i]|, по часовой стрелке начиная с верхнего левого угла самой метки. Это \emph{пиксели кадра}, а не метры сцены.
    \item \verb|rejected| --- отвергнутые кандидаты (квадраты, не подошедшие к словарю); для работы обычно не нужны.
\end{itemize}

В кадр может попасть несколько меток, и порядок в списке не определён. Нужную метку выбирают \emph{по значению ID}, а углы берут по тому же индексу:

\begin{codebox}[Выбрать метку 12]
\begin{verbatim}
if ids is not None:
    for i, value in enumerate(ids):     # i - номер в списке, value - [ID]
        if int(value[0]) == 12:
            points = corners[i][0]
\end{verbatim}
\end{codebox}

Функция \verb|enumerate| выдаёт на каждом шаге цикла пару: номер элемента и сам элемент.

\begin{taskbox}[Охота на маркеры]
Взлетите на 3~м нижней камерой симулятора, создайте детектор \texttt{DICT\_6X6\_250} и для одного кадра напечатайте «Вижу маркер X!» для каждого найденного ID. Затем создайте второй детектор с \texttt{DICT\_4X4\_50}, обработайте тот же кадр и напечатайте его список. Отлетите на $(0, 5, 3)$ и убедитесь, что программа сообщает «меток нет», не пытаясь читать ID.
\expected{над стартом --- строки \texttt{Вижу маркер 10!}, \texttt{Вижу маркер 11!}, \texttt{Вижу маркер 12!} (порядок может быть любым), у чужого словаря --- пустой список; в точке $(0, 5, 3)$ --- \texttt{меток нет} без ошибки.}
\mission{Python}{82--86}
\end{taskbox}

\subsection{Геометрия метки: углы, центр, ошибка}
\label{sec:python-aruco-geometry}

По четырём углам считают \textbf{центр} метки --- среднее координат: $c_x = (x_1 + x_2 + x_3 + x_4) / 4$, $c_y$ --- так же. Сравнивая центр метки с центром кадра, получают \textbf{ошибку} в пикселях:
$$ e_x = c_x - \frac{w}{2}, \qquad e_y = c_y - \frac{h}{2}. $$
Знак показывает сторону: $e_x > 0$ --- метка правее центра, $e_y > 0$ --- \emph{ниже} (ось $y$ кадра направлена вниз). Величина --- насколько далеко. Это ещё не команда дрону, а только измерение.

Одного кадра мало, чтобы поверить в находку: на шумном изображении детектор может ошибиться. Метку \textbf{подтверждают} несколькими кадрами подряд: счётчик растёт при каждом обнаружении и обнуляется при потере. Три обнаружения «с перерывом» --- не подтверждение.

\begin{codefile}[python]{python/aruco_examples/11_01_find_marker.py}
\end{codefile}

\paragraph{Что происходит.} Функция \verb|read_target| объединяет всё из этого урока: читает кадр, проверяет \verb|None|, запускает детектор, обрабатывает пустой результат и выбирает метку по ID. Основной цикл считает центр и ошибку, обводит метку линиями по углам (\verb|(k + 1) % 4| после последнего угла возвращает к первому), рисует центр и ведёт счётчик подряд идущих обнаружений. Дрон в этой программе только висит: измерение отделено от движения.

\begin{taskbox}[Ошибка относительно центра]
Измените пример: печатайте ошибку только знаками --- «метка правее/левее центра», «выше/ниже центра». Затем поставьте дрон в $(0.3, 0, 3)$ и в $(0.9, 0, 3)$ и сравните знак $e_x$. Подтверждение сделайте строже: пять кадров подряд, а на третьем шаге цикла отлетите на $(0, 5, 3)$ и вернитесь.
\expected{над стартом метка 12 правее и выше центра ($e_x > 0$, $e_y < 0$); из $(0.9, 0, 3)$ знак $e_x$ меняется на минус; после отлёта счётчик обнуляется, и «Подтверждено» печатается только после пяти новых обнаружений подряд.}
\mission{Python}{87--90}
\end{taskbox}

\subsection{ArUco на настоящем дроне}
\label{sec:python-aruco-real}

Дальше --- программы для компьютера, подключённого к Пионеру. В них два отличия от учебных: камера одна, \verb|Camera()| без \verb|view|, а распознавание идёт по настоящим пикселям, поэтому маркер печатают на бумаге или показывают на экране телефона.

\begin{codefile}[python]{python/aruco_examples/11_aruco_basics_example1_detect.py}
\end{codefile}

\subsubsection*{Разбор кода}
\begin{enumerate}
    \item \texttt{camera.get\_cv\_frame()} получает изображение с камеры дрона.
    \item \texttt{cv2.cvtColor} переводит картинку в оттенки серого (так алгоритм работает быстрее и надёжнее).
    \item \texttt{detect\_markers\_compat} вызывает детектор в форме, которую поддерживает установленная версия OpenCV: \texttt{ArucoDetector} (OpenCV 4.7 и новее) или старую функцию \texttt{aruco.detectMarkers}.
    \item Результат --- те же \texttt{corners} и \texttt{ids}, что в учебном примере; \texttt{drawDetectedMarkers} обводит найденные метки.
\end{enumerate}

Сама программа оформлена как \textbf{класс} (\texttt{class ArucoDetectorThread}): данные хранятся в полях \texttt{self.\dots}, а работа разложена по методам. Методы \texttt{capture\_thread} и \texttt{processing\_thread} запускаются в отдельных \textbf{потоках} (\texttt{threading.Thread}) и работают одновременно с показом кадра; общий кадр защищён замком \texttt{frame\_lock}, чтобы потоки не читали и не писали его в один момент. Потоки подробнее --- в разделе про групповой полёт этой главы. Пример использует словарь \texttt{DICT\_6X6\_50}: печатайте маркеры с номерами 0--49.

\subsubsection*{Определение координат маркера (только на дроне)}
Зная размер маркера и параметры камеры (полученные при калибровке --- см. раздел «Калибровка камеры» ниже), можно вычислить расстояние и ориентацию относительно маркера. В симуляторе \texttt{cv2.solvePnP} всегда возвращает неудачу, а файл калибровки не читается.

\begin{codefiletitle}[python]{Определение координат}{python/aruco_examples/11_aruco_basics_example2_coordinates.py}
\end{codefiletitle}

\subsubsection*{Посадка по маркеру (только на дроне)}
В этом примере дрон взлетает, занимает точку $(0, 0, 1)$ над стартом и висит, пока камера ищет маркеры. Как только в кадре появляется маркер с ID~42, дрон садится. Логика полёта работает в третьем потоке и только читает список \texttt{detected\_ids}, который заполняет поток обработки. Это посадка «увидел --- сел» без наведения: дрон сядет там, где висел. Как сначала навестись на метку, а потом сесть точно на неё, --- в следующем разделе.

\begin{codefiletitle}[python]{Полет по маркеру}{python/aruco_examples/11_aruco_basics_example3_flight.py}
\end{codefiletitle}

\begin{taskbox}[Маркер на бумаге]
Только на дроне. Распечатайте три маркера словаря \texttt{DICT\_6X6\_50} с разными номерами (или откройте их на экране телефона). Измените первый пример так, чтобы он печатал «Вижу маркер X!» один раз при появлении метки в кадре, а не на каждом кадре, и «Потерял маркер X» при её исчезновении.
\expected{при показе маркера в терминале одна строка о появлении, при уходе из кадра --- одна строка о потере; маркер из другого словаря программа не замечает.}
\end{taskbox}
